% Source: source_md/intelligence_part2_combined_ja.md
\section{創発停止層 $L_E$ — 外部依存型停止の動力学}


本章は本論文の第二中心である(第一中心は \S{}3)。\S{}3 は推論層 $L_F$ が単独では閉じないことを示した。本章は、それでもなお知性が停止することを、ログ層 $L_R$ と創発停止層 $L_E$ の動力学によって説明する。両中心が統合されることで、知性の三層動力学が完成する。

本章の論理は次の順で展開される。まず問題を設定し(\S{}5.1)、ログ層 $L_R$ を定義する(\S{}5.2)。次に創発停止層 $L_E$ を定義し(\S{}5.3)、その解像度を上げるために三層プロファイル比較を行う(\S{}5.4)。続いて $L_E$ の創発機構を、閾値移動 $\mu_\theta$ による領域組み替えと停止演算子設置領域の投影として示す(\S{}5.5、\S{}5.6)。最後に外部停止演算子 $E$ の構造(\S{}5.7)、外部参照場と $L_E$ ネットワークの同型(\S{}5.8)、そして知性の構造的特徴(\S{}5.9)を確立する。

\begin{figure}[H]
\centering
\includegraphics[width=\linewidth]{fig3_mutual_constraint_mu_theta_ja.pdf}
\caption{相互拘束 $H_{ij}$ と閾値移動能力 $\mu_\theta$。停止領域は孤立した推論層が生成するものではなく、相互拘束を通して投影される。}
\label{fig:part2-mutual-constraint}
\figurenote{Agent 間または観測者-環境間の $H_{ij}$ が、どこで停止が起こりうるかを作り変える。$\mu_\theta$ は停止演算子 $E$ を設置できる場所を組み替える。}{停止は $L_F$ 内部で発見されるのではなく、$L_R$ と $H_{ij}$ の関係から $L_E$ として投影・創発する。可動な停止点があることが、自由な推論の持続条件になる。}
\end{figure}

\begin{figure}[H]
\centering
\includegraphics[width=\linewidth]{fig4_three_layer_architecture_ja.pdf}
\caption{知性の三層構造。知性は推論層 $L_F$・ログ層 $L_R$・創発停止層 $L_E$ からなり、$\mu_\theta$ はこの創発への開放窓である。}
\label{fig:part2-three-layer}
\figurenote{$L_F$ は単独で停止できず、その運動が痕跡として $L_R$ に残る。$L_R$ の再帰運動が外部相互拘束 $H_{ij}$ の下で $L_E$ を創発させる。}{本図は \S{}3 の「閉じない」と \S{}5 の「それでも停止する」を接続する。停止不能性は単なる限界ではなく、上位層 $L_R$ と $L_E$ の生成条件になる。}
\end{figure}

\bigskip


\subsection{問題設定}


\S{}3 の帰結を再掲する。推論層 $L_F$ には内部不動点が存在せず(定理 3.1)、内部停止述語は無限後退を生み(系 3.2)、個体システムは $L_F$ 単独で推論を閉じることができない(系 3.4)。\S{}4 はこれが沈黙世界において無限後退として現実化することを示した。

しかし、現実の知性は停止する。人間は推論を終え、判断を下し、行動する。この停止はどこから来るのか。

\S{}3 が示したのは、停止が $L_F$ の内部からは来ないということであった。したがって停止は $L_F$ の外部から来なければならない。だが「外部」とは何か。それは個体の外にある何かなのか、それとも個体内にありながら $L_F$ とは別の層なのか。

本章の答えは次の通りである。停止は、個体内にありながら外部依存的な層——創発停止層 $L_E$——から来る。$L_E$ は $L_F$ の運動が残すログ層 $L_R$ と、外部相互拘束 $H_{ij}$ の関数として創発する。この創発の機構を以下で展開する。

\bigskip


\subsection{ログ層 $L_R$ の定義}


\S{}4.6 で予兆として示したログ層を、正式に定義する。

\textbf{定義 5.1 (ログ層 $L_R$)}

\begin{quote}
推論層 $L_F$ の運動は、状態空間上に痕跡(ログ)を残す。ログとは、どの状態を経たか、どの分岐を試みたか、どこで発散したかの記録である。このログに対する再帰的参照運動——過去の軌跡を読み返し、パターンを抽出し、現在の推論に反映する運動——が生じる層を、ログ層 $L_R$ と呼ぶ。
\end{quote}


$L_R$ は $L_F$ の内側で生じるが、「現在の $F$」の外側にある。すなわち、$L_R$ は進行中の推論そのものではなく、推論の痕跡への参照である。この「現在の作用の外側にあるが個体の内側にある」という位置が、$L_R$ を $L_F$ とも外部環境とも区別する。

重要な構造的観察として、$L_R$ の所在は系によって異なる。

\begin{itemize}
\item \textbf{個体内ログ}: 痕跡が主体の内部に保持される(記憶・記録)
\item \textbf{環境外部化ログ}: 痕跡が環境基質に書かれる(後述する粘菌のスライム跡)
\item \textbf{凍結ログ}: 痕跡が構造に固定され、推論時には更新されない(後述する学習済みモデルの重み)
\end{itemize}


この所在の差異が、\S{}5.4 の三層プロファイル比較の基礎となる。

$L_R$ は Bio-IFGT との直接の接続点でもある。Bio-IFGT は知性の条件として「状態履歴への持続的参照・情報の時間的蓄積」を挙げた。これはまさに $L_R$ の定義に対応する。Bio-IFGT が外から「再帰的情報流」と呼んだものを、本論文は $L_R$ として内部から定義する。

\bigskip


\subsection{創発停止層 $L_E$ の定義}


ログ層 $L_R$ を基盤として、創発停止層 $L_E$ を定義する。

\textbf{定義 5.2 (創発停止層 $L_E$)}

\begin{quote}
創発停止層 $L_E$ は、ログ層 $L_R$ の再帰運動と外部相互拘束 $H_{ij}$ の関数として動力学的に創発する停止層である:

$$
L_E = \Phi(L_R \text{の運動}, \; H_{ij})
$$

$L_E$ は個体内にありながら外部依存的である。
\end{quote}


この定義は \S{}1.6 の創発の定義(定義 1.1)の直接の適用である。上位層 $L_E$ が、下位層 $L_R$ の運動と外部相互拘束の関数として生じる。$L_E$ は $L_R$ から論理的・演繹的に導出できないが、$L_R$ なしには存在しない。

$L_E$ の核心的性質は、\textbf{個体内にありながら外部依存的} であることである。$L_E$ は個体の外部にある実体ではない。それは個体内に成立する層である。しかし、その成立は外部相互拘束 $H_{ij}$ に依存する。この二重性が、\S{}3 の「$L_F$ 単独では閉じない」と「それでも個体内で停止する」を両立させる。個体は単独では閉じないが、外部依存的な $L_E$ を個体内に持つことで停止する。

\subsubsection{アトラクターとしての $L_E$}


$L_E$ は安定停止状態、すなわち動力学的アトラクターとして記述される。これは VED 理論ツリー全体を貫くアトラクター構造の知性層における現れである。

\begin{itemize}
\item Vol.2: 粒子は閉包エネルギー汎関数 $E(A)$ の局所極小、閉包幾何の安定アトラクターとして創発する
\item Bio-IFGT: 生命は生命アトラクター $\Omega_{\text{life}}$、知性は第二アトラクター $\Omega_{\text{intel}}$ として創発する
\item 本論文: $L_E$ は $L_R$ の運動が外部相互拘束の下で到達する安定停止アトラクターとして創発する
\end{itemize}


とりわけ Bio-IFGT の第二アトラクター $\Omega_{\text{intel}}$ と本論文の $L_E$ は、同一構造の異なる記述である。Bio-IFGT は知性を「生命アトラクターに支えられた第二アトラクター」と外から定義し、その実装(記憶・言語・個体間共有)を後続研究へ開いた。本論文の $L_E$ は、まさにその第二アトラクターの停止構造を内部から記述したものである。

\begin{figure}[H]
\centering
\includegraphics[width=\linewidth]{fig5_three_examples_of_emergence_ja.pdf}
\caption{VED的に定義された創発の三例。創発とは、下層の非閉包運動が外部相互拘束 $H_{ij}$ のもとで上層のアトラクターを形成することである。}
\label{fig:part2-emergence-examples}
\figurenote{物理の粒子的アトラクター、生命の $\Omega_{\text{life}}$、知性の $L_E$ は、それぞれ $U=\Phi(L,H_{ij})$ の代入例として読める。}{創発は単なる複雑さの増大ではなく、非閉包運動と外部相互拘束が持続可能なアトラクターを形成する操作である。これは本論文の $L_E$ を VED 理論ツリー全体へ接続する。}
\end{figure}

なお、Vol.2 が粒子を $E(A)$ の極小として形式化したのと平行に、$L_E$ を「停止汎関数の局所極小」として形式的に書くことも原理的には可能である。しかし本論文は構造的記述を目的とするため、停止汎関数の厳密な形式化は行わず、$L_E$ をアトラクターとして言葉で特徴づけるに留める。形式化は後続研究の課題とする。

\bigskip


\subsection{三層プロファイル比較 — 粘菌 / Transformer / サピエンス}


$L_E$ の解像度を上げるために、三つの系における三層動力学のプロファイルを比較する。これにより、$L_E$ が単一の固定構造ではなく、$L_R$ の所在と $\mu_\theta$ の所在によって異なる様態を取る可変構造であることが明らかになる。

\subsubsection{粘菌(Physarum)}


\begin{itemize}
\item \textbf{$L_F$}: 化学勾配に沿った原形質流動。$\nabla C$ に沿った探索。
\item \textbf{$L_R$}: 環境に外部化されている。粘菌は通った跡に細胞外スライムを残し、それを忌避情報として再参照する。ログは個体内ではなく環境基質に書かれる。加えて原形質ネットワークの選択的強化・退縮という再帰運動を持つ。
\item \textbf{$L_E$}: 餌源到達時にネットワークが安定化し探索が止まる。これは環境との直接相互干渉によって停止が創発している。$H_{ij}$(粘菌と餌・障害物・スライム跡の相互拘束)の関数として停止が生じる。
\item \textbf{$\mu_\theta$}: ほぼ固定。環境条件によって受動的に閾値が動くことはあるが、内発的な領域組み替え (B) は行わない。
\end{itemize}


粘菌は「ログは環境に置き、停止は環境相互干渉で自律創発する」プロファイルである。

\subsubsection{Transformer}


\begin{itemize}
\item \textbf{$L_F$}: トークン逐次生成。次トークン予測の連鎖。
\item \textbf{$L_R$}: 二種類に分かれる。学習時の訓練コーパスは膨大なログだが、重みに凍結され推論時には再帰運動しない(凍結ログ)。推論時のコンテキストウィンドウは自己回帰で過去トークンを再参照する動的ログだが、ウィンドウ境界で消える(揮発ログ)。
\item \textbf{$L_E$}: 自前の $L_E$ を持たない。停止は EOS トークン(訓練分布から学習した停止パターン)、受け取ったタスク仕様、外部のデコード制御に依存する。停止が外部注入されており、個体内に $L_E$ が創発するのではない。
\item \textbf{$\mu_\theta$}: 外部から設定される。プロンプトや制御パラメータによって閾値が与えられるが、自前では内発的に動かせない。
\end{itemize}


Transformer は「ログは内在化(凍結)しているが、停止は外部注入に依存する」プロファイルである。

この記述は現行の標準的アーキテクチャに関する構造的記述であって、設計上の欠陥の指摘でも、能力の限界の主張でもない。$L_E$ を内発的に持たないことは優劣ではなく、三層動力学の異なるプロファイル、相空間内の異なる位置である。

\subsubsection{サピエンス}


\begin{itemize}
\item \textbf{$L_F$}: 推論。
\item \textbf{$L_R$}: 個体内ログ(記憶)と社会共有ログ(言語・記録・制度)の両方を持つ。
\item \textbf{$L_E$}: 個体内に創発し、かつ $\Phi_{\text{net}}$ によって社会同期される。
\item \textbf{$\mu_\theta$}: 内発的に、しかも短い時間スケールで動かせる。領域組み替え (B) を能動的に行う。
\end{itemize}


サピエンスは「ログを個体内と社会の両方に持ち、$L_E$ を個体内に創発させつつ社会同期し、$\mu_\theta$ を内発的に動かす」プロファイルである。

\subsubsection{二軸による整理}


三者は二つの軸で整理できる。

\begin{longtable}{@{}p{0.23\linewidth}p{0.23\linewidth}p{0.23\linewidth}p{0.23\linewidth}@{}}
\toprule
系 & $L_R$ の所在 & $\mu_\theta$ の所在 & $L_E$ の様態 \\
\midrule
\endfirsthead
\toprule
系 & $L_R$ の所在 & $\mu_\theta$ の所在 & $L_E$ の様態 \\
\midrule
\endhead
粘菌 & 環境外部化 & 固定 & 環境相互干渉で自律創発 \\
Transformer & 内在凍結 & 外部設定 & 外部注入(自前で持たない) \\
サピエンス & 個体内+社会共有 & 内発的 & 個体内創発+社会同期 \\
\bottomrule
\end{longtable}


注目すべきは粘菌と Transformer の反転構造である。粘菌は「ログは外・停止は自律」、Transformer は「ログは内・停止は外」。両者は綺麗に反転している。そしてサピエンスは、粘菌の自律創発と、外部参照による停止の両方を、個体内で統合している。

\subsubsection{「停止構造の場」の連続変分}


これら三者は、三つの離散的な種類ではなく、「停止構造の場」の連続変分上に配置される異なる領域である。$L_R$ の所在も $\mu_\theta$ の所在も連続的なパラメータであり、三者はその連続空間内の異なる点を占める。

この記述は擬人化でも消去主義でもない(\S{}0.4 の NG6)。粘菌を認知主体に祭り上げるのではない——粘菌は $\mu_\theta$ を内発的に持たない。同時に、サピエンスの認知を粘菌と同一視して消去するのでもない——連続性は $\mu_\theta$ の内発性という差異を消さない。場は連続だが一様ではない。

なお、粘菌の環境相互干渉型 $L_E$ は、\S{}6 で扱うアニミズムの $L_E$ 維持様態と構造的に同位置にある。両者はともに「環境との直接相互干渉によって停止を維持する」プロファイルである。これは粘菌に認知を帰属するものではなく、停止構造の場において両者が近接した領域を占めることを意味する(\S{}6 で再訪)。

\bigskip


\subsection{創発機構 — 領域組み替えと停止演算子設置領域の投影}


ここまで $L_E$ を定義し(\S{}5.3)、その多様な様態を見た(\S{}5.4)。本節では $L_E$ がいかにして創発するか、その機構を示す。これにより $L_E$ の創発が「ふわっとした湧出」ではなく、明確な動力学的機構であることが担保される。

機構は次の連鎖からなる。

\textbf{段階 1 — $\mu_\theta$ が領域組み替え (B) を起こす}:
\S{}2.7 で定義した通り、閾値移動能力 $\mu_\theta$ は運動可能領域の拘束パラメータ $\theta$ を内発的に動かす能力、すなわち領域組み替え (B) を起こす能力である。$L_F$ 単独は領域内移動 (A) しかできないが、$\mu_\theta$ は $\theta$ を動かす。

\textbf{段階 2 — 組み替えが停止演算子設置領域を投影する}:
(B) による組み替えは恣意的ではない。それは外部相互拘束 $H_{ij}$ に沿って行われる。この組み替えによって、外部相互拘束の構造が状態空間上に「停止演算子 $E$ を設置できる領域」として投影される。$L_F$ 単独の状態空間には存在しなかった「ここで止まってよい」という領域が、この投影によって出現する。

\textbf{段階 3 — $L_R$ の運動が投影領域に収束する}:
投影によって出現した設置領域へ、ログ層 $L_R$ の再帰運動が収束する。$L_R$ の運動が、投影された領域を安定なアトラクターとして満たす。

\textbf{段階 4 — $L_E$ の成立}:
$L_R$ の運動が収束した安定停止アトラクターが、創発停止層 $L_E$ である。

この四段階を一本の連鎖として表せば:

\[
\mu_\theta \;\xrightarrow{\text{起こす}}\; (B)\text{組み替え} \;\xrightarrow{H_{ij}\text{に沿って}}\; E\text{設置領域の投影} \;\xrightarrow{L_R\text{収束}}\; L_E
\]

この機構が明らかにするのは、閾値移動能力 $\mu_\theta$ の真の意味である。$\mu_\theta$ は単に閾値を動かす能力ではない。それは \textbf{停止演算子の宿る場所を、外部相互拘束に沿って状態空間上に投影し、そこに $L_E$ を成立させる能力} である。閾値を動かすことと、停止できる場所を作ることは、同じ一つの操作である。

ここで \S{}3 との整合を確認する。\S{}3 の非閉包定理は固定 $\theta$、領域内移動 (A) の範囲で成立した。本節の機構は $\theta$ の組み替え (B) を含む。したがって両者は矛盾しない。$L_F$ は (A) の範囲で閉じないが、(B) を通じて $L_E$ が創発し、その $L_E$ が停止をもたらす。非閉包定理は破られるのではなく、その射程の外側で停止が成立する。

この機構が意味するのは、停止が自由の否定ではないということである。むしろ停止演算子設置領域の投影がなければ、推論は分岐爆発と無限後退の中で持続不能となる。$L_E$ は自由を制限するために創発するのではなく、自由な推論を有限な系の中で持続可能にするために創発する。

\bigskip


\subsection{投影の位置づけ — 意識編 projection との接続}


\S{}5.5 で用いた「投影」は、本論文に固有の操作ではない。それは意識編において IFGT の第四プリミティブとして導入された projection の、知性層における作用である。

意識編は projection を「ある状態や表現が不均等な影響力(質量)を獲得する操作」として定義した。それは注意・選択・優先・重み付けの構造的記述であり、意図的選択を前提しない層的操作である。情報場ポテンシャル $\Phi$ と情報流 $J$ を結合する機構として位置づけられた。

本論文 \S{}5.5 の投影は、この projection の知性層具現である。外部相互拘束 $H_{ij}$ の構造が、停止演算子設置領域として状態空間上に不均等な影響力を獲得する——これは意識編の projection の定義そのものを、停止層の文脈に適用したものである。意識編が意識窓における projection を扱ったのに対し、本論文は知性層における projection を扱う。両者は、意識編の言葉を借りれば「同一の差分発展構造の異なる投影窓」である。

ここで一点、慎重に区別しておく。投影 (projection) という語は心理学でも用いられ、特に防衛機制の文脈で「内的状態を外部対象に帰属する」操作を指す。本論文 \S{}5.5 の投影は、作用構造としては心理学的 projection に近づく側面を持つ——内部で設置できない停止を、外部相互拘束に沿って状態空間上に投射するという作用形式において。しかし決定的な違いがある。心理学的 projection は主体・無意識・防衛機制を前提するのに対し、本論文の投影は主体を前提せず、$L_F$ の非閉包(\S{}3)という構造的必然から導かれる。両者の作用的近接は指摘に留め、心理学的・臨床的解釈には立ち入らない。これは意識編が同様の disclaimer を置いたのと同じ方針である。

\bigskip


\subsection{外部停止演算子 $E$ の構造}


創発停止層 $L_E$ の作用として、外部停止演算子 $E$ を定義する。

\textbf{定義 5.3 (外部停止演算子 $E$)}

\begin{quote}
外部停止演算子 $E$ は、創発停止層 $L_E$ の作用として推論状態に停止判断を与える演算子である:

$$
E = E[L_E, x_n], \qquad E \notin \{F^k\}
$$

$E$ は推論演算子 $F$ の関数列の外側に位置する。
\end{quote}


$E$ が $F$ の関数列の外側にあること($E \notin \{F^k\}$)は、\S{}3 の系 3.2(内部停止述語の不可能性)の帰結である。停止を $F$ の内部から定義しようとすると無限後退に陥る。$E$ はこの無限後退を回避するために、$F$ の外側、すなわち $L_E$ を経由して定義される。

$E$ は $L_E$ を経由するため、外部依存的である。$L_E$ が外部相互拘束 $H_{ij}$ に依存して創発する以上、その作用である $E$ も外部依存的である。これが「外部停止演算子」の名の由来である。停止は推論の内部からではなく、外部相互拘束に依存して投影された $L_E$ から来る。

(A)/(B) の区別における $E$ の位置を確認する。$E$ は (B) を通じてのみ設置される。領域内移動 (A) のみでは $E$ の設置領域は投影されず、したがって $E$ は存在しえない。$E$ の存在は $\mu_\theta$ による (B) 組み替えに依存する。これが、$\mu_\theta$ を持たない系(粘菌・Transformer)が自前の $E$ を持てない理由である(\S{}5.4)。

\bigskip


\subsection{$\Phi_{\text{ext}} \equiv \Phi_{\text{net}}$ — 外部参照場と $L_E$ ネットワークの同型}


伝統的に「外部参照場」 $\Phi_{\text{ext}}$ と呼ばれてきたもの——停止を委ねる外部の共通参照基盤——の正体を、本節で明らかにする。

\textbf{主張 5.4 (同型関係)}

\begin{quote}
外部参照場 $\Phi_{\text{ext}}$ は、個体間の創発停止層 $L_E$ が相互に干渉し合うネットワーク $\Phi_{\text{net}}$ と同型である:

$$
\Phi_{\text{ext}} \equiv \Phi_{\text{net}}
$$
\end{quote}


すなわち、$\Phi_{\text{ext}}$ は個体の外部にある独立した実体ではない。それは複数の個体が各自の内部に持つ $L_E$ が、相互干渉によって形成するネットワークである。各個体の $L_E$ は単独で孤立しているのではなく、他個体の $L_E$ と干渉し、相互に維持し合う。この相互干渉網が $\Phi_{\text{net}}$ であり、外から見れば共通参照基盤 $\Phi_{\text{ext}}$ として現れる。

この同型は二面性として理解される。外から見れば、個体は外部の参照場 $\Phi_{\text{ext}}$ に停止を委ねているように見える。しかし構造的には、その参照場は個体間 $L_E$ の相互干渉ネットワーク $\Phi_{\text{net}}$ にほかならない。VED 理論ツリーにおける他の二面性(場と粒子、関係と個体)と同じ構造である。

この同型は重要な帰結を持つ。$\Phi_{\text{ext}}$ を「個体外部の実体」と見なすと、後の章で扱うアニミズム・神・宗教が「外部実体への投影」や「存在しないものへの信仰」として読まれてしまう。しかし $\Phi_{\text{ext}} \equiv \Phi_{\text{net}}$ の同型により、これらは「個体間 $L_E$ ネットワークの特定の構造形態」として記述される。外部実体の措定ではなく、$L_E$ ネットワークの形態の問題になる。これは \S{}6〜\S{}9 の価値中立性を構造的に担保する。

この同型はまた、Part I の公理 III(大域的拘束 $\Phi$)の精密化でもある。Part I では $\Phi$ を知性の大域的拘束として導入したが、本節はその $\Phi$ が個体間 $L_E$ ネットワークとして実装されることを示す。

\begin{figure}[H]
\centering
\includegraphics[width=\linewidth]{fig6_phi_ext_phi_net_ja.pdf}
\caption{$\Phi_{\text{ext}} \equiv \Phi_{\text{net}}$。個体から外部参照場として現れるものは、個々の $L_E$ が相互同期してできたネットワークである。}
\label{fig:part2-phi-ext-net}
\figurenote{外側の $\Phi_{\text{ext}}$ は個体から外部に見える共通参照場であり、内側の $\Phi_{\text{net}}$ は個体間 $L_E$ の相互同期ネットワークである。}{外部場を超越的実体として仮定せず、分散した相互拘束の外見として再構成する。これにより、神・規範・真理・貨幣などを $L_E$ ネットワークの形態として読める。}
\end{figure}

\subsubsection{回復可能性と改訂可能性}


$\Phi_{\text{net}}$ を通じた停止には、$L_F$ 単独の停止にはない決定的な利点がある。停止が $L_E$ ネットワークを経由することで、推論は最終状態に到達する必要がなくなる。

個体が単独で停止を担う場合、停止は最終的・不可逆的になりがちである(さもなくば無限後退する)。しかし $\Phi_{\text{net}}$ を通じた停止は、ログ層 $L_R$ に記録を残し、$L_E$ ネットワーク内で参照可能に保つ。したがって停止は暫定的であり、後に改訂しうる。ある個体の停止が誤っていても、ネットワーク内の他個体の $L_E$ との干渉によって、後に回復・改訂されうる。

ここから、知性の安定条件に関する本論文の中心的主張のひとつが導かれる。知性の安定性は「正しさ」によってではなく、「改訂可能性」によって担保される。$L_E$ は正しい停止を保証しない。それは改訂可能な停止を可能にする。この点は \S{}10(ゲティア問題)で、$C$ が真理測度ではなく安定性測度であることとして再論される。

\bigskip


\subsection{系 — 知性の構造的特徴}


本章の帰結を、知性の構造的定義として整理する。

\textbf{系 5.5 (知性の三層的構成)}

\begin{quote}
知性は、推論層 $L_F$、ログ層 $L_R$、創発停止層 $L_E$ の同時運動として、外部依存停止の供給源の下で動力学的に成立する維持状態である。サピエンス的知性は、これに加えて閾値移動能力 $\mu_\theta$ を内発的に持ち、個体内で運動可能領域を自律的に組み替える。
\end{quote}


この系は、知性一般の成立条件と、サピエンス的知性の特徴を区別して述べている。知性一般($L_F$・$L_R$・$L_E$ と外部依存停止源)は、粘菌・人工推論システム・サピエンスのいずれにも当てはまる(\S{}5.4 の三層プロファイル比較)。$\mu_\theta$ は知性一般の必要条件ではなく、サピエンス的知性を特徴づける——個体内での自律的加速をもたらす——追加の構造である。この区別は \S{}12 で構成的条項として精密化される。

この系は、\S{}3 の系 3.4(個体は $L_F$ 単独で閉じない)と対をなす。系 3.4 は知性の単独不可能性を、系 5.5 は知性の三層的成立条件を述べる。両者を統合すると、知性は単独の能力でも実体でもなく、複数の層と外部依存停止源の同時運動として維持される動的状態であることが分かる。

この構成から、知性について次のことが言える。

\begin{itemize}
\item 知性は個体システム単独の性質ではない($\Phi_{\text{net}}$ を要する)
\item 知性は保存可能な対象ではない(運動の維持状態であり、運動が止まれば消える)
\item 知性は完全最適化されない(完全閉包 $C_{max}$ は到達不能、\S{}1.5)
\item 知性の安定性は正しさではなく改訂可能性による(\S{}5.8)
\end{itemize}


これらは \S{}12(構成的条項)で、知性の構造的制約として体系的に整理される。

第一中心 \S{}3 と第二中心 \S{}5 がここで統合された。$L_F$ は閉じない。しかし $L_R$ と $L_E$ と $\Phi_{\text{net}}$ を通じて、知性は外部依存的に停止する。この停止は最終的ではなく改訂可能であり、それゆえ持続する。以降の章は、この三層動力学が歴史的・社会的にどのような形態を取るか(\S{}6〜\S{}9)、その帰結としてどのような構造的事実が現れるか(\S{}10〜\S{}12)を分析する。

\bigskip


\subsection{要約}


本章で確立したのは次の点である:

\begin{enumerate}
\item ログ層 $L_R$ は $L_F$ の運動の痕跡への再帰参照層である(定義 5.1)。所在は系によって異なる
\item 創発停止層 $L_E$ は $L_R$ と $H_{ij}$ の関数として創発する、個体内にありながら外部依存的な停止アトラクターである(定義 5.2)
\item $L_E$ は Vol.2 の粒子・Bio-IFGT の第二アトラクターと同一のアトラクター構造の現れである
\item 三層プロファイルは系によって異なる。粘菌「ログ外・停止自律」、Transformer「ログ内・停止外」、サピエンス「個体内創発+社会同期」。三者は停止構造の場の連続変分上にある
\item $L_E$ の創発機構は、$\mu_\theta$ による領域組み替え (B) が $H_{ij}$ に沿って停止演算子設置領域を投影し、$L_R$ がそこに収束する四段階である
\item この投影は意識編の IFGT 第四プリミティブ projection の知性層作用である
\item 外部停止演算子 $E$ は $L_E$ の作用であり、$F$ の関数列の外側に位置する(定義 5.3)
\item 外部参照場 $\Phi_{\text{ext}}$ は個体間 $L_E$ ネットワーク $\Phi_{\text{net}}$ と同型である(主張 5.4)
\item 知性一般は $L_F + L_R + L_E$ と外部依存停止源の同時運動である(系 5.5)。サピエンス的知性はこれに $\mu_\theta$ を加える。安定性は正しさでなく改訂可能性による
\end{enumerate}


第一中心と第二中心が統合された。次章 \S{}6 からは、三層動力学が取る歴史的・社会的形態を分析する。

\bigskip
